\section{Overview}
\label{sec:overview}

RAFTcorr measures full-field displacement and strain from a sequence of images of a deforming
specimen. The first image of the sequence is the reference image. Every other image is compared
with it, or with a reference frame earlier in the sequence, by a learned correlation model: a neural network trained
on speckle images whose displacement is known exactly. The model returns one displacement vector
for every pixel of the region of interest, so there is no subset size and no mesh to choose.
Strain follows from the displacement with a virtual strain gauge, in the strain window.

The same engine runs behind three interfaces: the desktop application, which most of this guide
follows; the \file{raftcorr} command line; and the Python library (Section~\ref{sec:cli}). A
project saved by one of them opens in the others.

\begin{notebox}
RAFTcorr is two-dimensional digital image correlation: one camera, a planar specimen, motion in
the image plane. It does not reconstruct out-of-plane motion.
\end{notebox}

\subsection{The four models}
\label{sec:overview-models}

RAFTcorr ships four models, \file{RAFTcorr-s1}, \file{-s2}, \file{-s4} and \file{-s8}. They share
one network architecture and differ in one number, the encoder downsampling factor $s$: the
network describes both images on a grid $s$ times coarser than the pixels, and the correlation
search at its coarsest level looks $32\,s$ px away. That distance is the model's nominal
\emph{reach}.

\begin{table}[htbp]
  \centering
  \caption{The four models. The \ui{Model} list shows each with its reach, for example
  \ui{RAFTcorr-s4 (128 px)}.}
  \label{tab:models}
  \begin{tabular}{lccl}
    \toprule
    Model & $s$ & Reach ($32\,s$) & Use it for \\
    \midrule
    RAFTcorr-s1 & 1 & 32 px  & small displacements, finest detail (strain concentrations) \\
    RAFTcorr-s2 & 2 & 64 px  & small to moderate displacements, fine detail, faster than s1 \\
    RAFTcorr-s4 & 4 & 128 px & a balance of detail and reach (the default) \\
    RAFTcorr-s8 & 8 & 256 px & large motion and rotation, or no prior knowledge \\
    \bottomrule
  \end{tabular}
\end{table}

A smaller $s$ resolves shorter wavelengths of the displacement; a larger $s$ reaches further. The
reach has to cover the largest displacement of any image pair, so the choice depends on the
tracking mode too (Section~\ref{sec:workflow-model}). When you do not know how far the specimen
moves, \ui{Suggest a Model} measures the motion of your sequence with s8 and picks the finest
model whose reach covers it.

Each model file is about 20\,MB. RAFTcorr downloads a model the first time an analysis uses it
and keeps it for later runs (Section~\ref{sec:installing-models}).

\subsection{An analysis in eight steps}
\label{sec:overview-steps}

The main window has three columns: settings on the left, the image in the middle, the run controls
and the display settings on the right. An analysis goes through it from left to right.

\begin{enumerate}
  \item \textbf{Load the images} in the \ui{IMAGES} section: a folder of images, sorted in natural
    order. The first image is the reference image (Section~\ref{sec:images}).
  \item \textbf{Choose the tracking mode and the model} in \ui{WORKFLOW TYPE}. Do this before the
    region of interest: the tracking mode decides which frames are reference frames
    (Section~\ref{sec:workflow}).
  \item \textbf{Define the region of interest} on frame~1 in \ui{REGION OF INTEREST}, and on
    later frames that must leave a region out, such as a crack or a hole that opens
    (Section~\ref{sec:roi}).
  \item \textbf{Check the parameters} in \ui{PARAMETERS} and \ui{ADVANCED}. The defaults suit most
    sequences (Section~\ref{sec:parameters}).
  \item \textbf{Run} with \ui{Run DIC Analysis}, and follow the run in \ui{PROGRESS} and \ui{LOG}
    (Section~\ref{sec:running}).
  \item \textbf{Review the displacement} on the canvas: each component, the magnitude, and the
    residual, which shows how well each frame was tracked (Sections~\ref{sec:canvas}
    and~\ref{sec:results}).
  \item \textbf{Compute the strain} in the strain window, which opens when the run finishes, and
    measure points, lines and areas in its \ui{Analysis} tab (Sections~\ref{sec:strain}
    and~\ref{sec:analysis}).
  \item \textbf{Export and save}: data files, figures, animations and a report with
    \ui{Export Results}, and the whole session with \menu{File}{Save Project}
    (Sections~\ref{sec:export} and~\ref{sec:projects}).
\end{enumerate}

\screenshot{fig-main-window.png}{The main window after a run, showing the vertical displacement on
the last frame. The numbers mark the five regions described in the text: (1) the \ui{IMAGES}
section; (2) the settings column; (3) the canvas with the frame navigator; (4) the run controls and
\ui{PROGRESS}; (5) \ui{FIELD}, \ui{VISUALIZATION}, \ui{PHYSICAL UNITS} and \ui{LOG}.}{fig:main-window}

The five regions of Figure~\ref{fig:main-window}:

\begin{description}[leftmargin=0.6cm,style=nextline]
  \item[1. IMAGES]
    The drop zone (\ui{Drop image folder or Browse}), the \ui{Natural Sort (1, 2, …, 10)}
    checkbox and the image list. The list numbers the frames from 1 and has a \ui{Region}
    column, whose button adds or edits the region of interest of each frame. See
    Section~\ref{sec:images}.
  \item[2. The settings column]
    Four sections that fold open and closed: \ui{WORKFLOW TYPE} (tracking mode, model,
    \ui{Suggest a Model}), \ui{REGION OF INTEREST} (drawing tools and mask files),
    \ui{PARAMETERS} (iterations, smoothing, device) and \ui{ADVANCED} (margin, GPU memory, tile
    blending, correlation; folded by default). See Sections~\ref{sec:workflow},
    \ref{sec:roi} and~\ref{sec:parameters}.
  \item[3. The canvas]
    The current frame with the region of interest or the displacement field drawn over it, the
    zoom buttons and a read-out of the position and value under the cursor at the top, and the
    frame navigator at the bottom. See Section~\ref{sec:canvas}.
  \item[4. Run controls and PROGRESS]
    \ui{Run DIC Analysis}, \ui{Cancel}, \ui{Export Results} and \ui{Open Strain Window}, and below
    them the progress of a run. See Sections~\ref{sec:running}, \ref{sec:strain}
    and~\ref{sec:export}.
  \item[5. Display and log]
    \ui{FIELD} chooses what is drawn (\ui{Disp U}, \ui{Disp V}, \ui{Magnitude}, \ui{Residual}) and
    where (\ui{Show on}); \ui{VISUALIZATION} holds the colormap, the colour range and the opacity;
    \ui{PHYSICAL UNITS} converts pixels and frames to lengths and seconds; \ui{LOG} lists every
    message of the session with its time, warnings in yellow and errors in red. See Section~\ref{sec:results}.
\end{description}

The menu bar has three menus: \ui{File} (images, projects, export), \ui{Settings} (theme and
language) and \ui{Help} (models, self-test, version). Section~\ref{sec:settings} lists them with
their keyboard shortcuts.

\subsection{Conventions}
\label{sec:conventions}

These rules hold everywhere: on the screen, in exported files, on the command line and in
Python.

\begin{itemize}
  \item \textbf{Displacement} $(u, v)$ goes from the reference image to the deformed image: a
    material point at pixel $X$ in the reference image appears at $X + (u, v)$ in the deformed
    image. $+u$ points right and $+v$ points down, as the image rows and columns run. The unit is
    the pixel, unless you set physical units (Section~\ref{sec:physical-units}).
  \item \textbf{Lagrangian storage.} Every field is stored in the coordinates of the reference
    image: the value at a pixel belongs to the material point that sits at that pixel in frame~1.
    Incremental tracking composes its steps into these coordinates as well. Choosing
    \ui{Deformed frame} in \ui{Show on} moves the field to the current positions of the points for
    display only; the stored data do not change.
  \item \textbf{Frame numbers.} The application counts frames from 1: the image list, the frame
    navigator and the log call the reference image frame~1. File names written by
    the export dialog count the same way (\file{frame\_0001} is the reference image). Inside the
    arrays of the exported data files, and in Python, index $k$ is image $k$ counted from 0, so
    frame~$n$ on the screen is index $n-1$. The command line uses both: its progress lines, its
    messages and its \file{suggest} report count from 1, its frame option \file{-{}-keys} from 0
    (Section~\ref{sec:cli}).
  \item \textbf{Shear strain} $\varepsilon_{xy}$ is the tensor component, half the engineering
    shear strain $\gamma_{xy}$ (Section~\ref{sec:strain}).
  \item \textbf{Grey levels.} RAFTcorr compares grey levels from 0 to 255. Colour images are
    converted to luminance, and images with more than 8 bits are scaled with one factor for the
    whole sequence, taken from the reference image (Section~\ref{sec:images}).
\end{itemize}

\subsection{Where RAFTcorr keeps its files}
\label{sec:overview-files}

RAFTcorr writes into the image folder only what you save there yourself. Its own files are in
your user account, so no administrator rights are needed.

\begin{description}[leftmargin=0.6cm,style=nextline]
  \item[Model files]
    \file{\%LOCALAPPDATA\%\textbackslash RAFTcorr\textbackslash models} on Windows,
    \file{\textasciitilde/.raftcorr/models} on macOS and Linux. \menu{Help}{Models…} shows the
    folder and opens it (Section~\ref{sec:installing-models}).
  \item[Settings]
    The theme, the language and the last export folder, in the settings store of your user
    account: the registry key
    \file{HKEY\_CURRENT\_USER\textbackslash Software\textbackslash RAFTcorr} on Windows, a property
    list in \file{\textasciitilde/Library/Preferences} on macOS.
  \item[Log file]
    \file{RAFTcorr.log}, however RAFTcorr was installed, in this folder:\newline
    Windows: \file{\%LOCALAPPDATA\%\textbackslash RAFTcorr\textbackslash logs}\newline
    macOS: \file{\textasciitilde/Library/Application Support/RAFTcorr/logs}\newline
    Linux: \file{\$XDG\_DATA\_HOME/RAFTcorr/logs}, by default
    \file{\textasciitilde/.local/share/RAFTcorr/logs}\newline
    At 2\,MB the
    file is renamed and a new one started; the three most recent older files are kept
    (Section~\ref{sec:troubleshooting}).
  \item[Windows installation]
    \file{\%LOCALAPPDATA\%\textbackslash Programs\textbackslash RAFTcorr} by default, with the
    private Python and \file{install.log}, the log of the setup (Section~\ref{sec:installing}).
  \item[Large results]
    When the results of a run (displacement and residual) need more than 2\,GB of memory,
    RAFTcorr keeps them in memory-mapped files in a temporary folder named
    \file{raftcorr\_results\_\ldots}, which the \ui{LOG} names, and deletes the folder when you
    quit. Save a project or export the results to keep them (Sections~\ref{sec:projects}
    and~\ref{sec:export}).
\end{description}
